Le chlore existe dans la nature sous forme de trois isotopes~: le chlore 35
$(\ce{_{17}^{35}Cl})$, le chlore 36 $(\ce{_{17}^{36}Cl})$ et le chlore 37
$(\ce{_{17}^{37}Cl})$.

\begin{donnees}
    \begin{table}[H]
        \centering
        \setlength{\arrayrulewidth}{0.8pt}
        \renewcommand{\arraystretch}{1.3}
        \begin{tabularx}{\textwidth}{|p{4cm}|C|C|C|p{4cm}|}
            \hline
            Noyau ou particule & Électron & Chlore $\ce{_{17}^{36}Cl}$ &
            Argon $\ce{_{18}^{36}Ar}$ &
            \multirow{2}{*}{$\qty{1}{\u}=\qty{931.5}{\U}$} \\
            \cline{1-4}
            Masse en (\unit{\u}) & $0,000549$ & $35,968312$ & $35,967545$ & \\
            \hline
        \end{tabularx}
        \begin{tabularx}{\textwidth}{|p{10cm}|C|C|C|}
            \hline
            Noyau & $\ce{_{17}^{35}Cl}$ & $\ce{_{17}^{36}Cl}$ & 
            $\ce{_{17}^{37}Cl}$ \\
            \hline
            \raisebox{0mm}[7mm][4mm]{
                Énergie de liaison par nucléon~:
                $\frac{E_{\ell}}{A}$ ($\unit{\MeV\per\mathrm{nucléon}}$)} &
            $8,5178$ & $8,5196$ & $8,5680$ \\
            \hline
        \end{tabularx}
    \end{table}
\end{donnees}

\begin{enumerate}
    \item Déterminer, en justifiant votre réponse, le noyau le plus stable parmi
          $\ce{_{17}^{35}Cl}$, $\ce{_{17}^{36}Cl}$ et
          $\ce{_{17}^{37}Cl}$.~\textbf{(1pt)}

    \item Le chlore 36 radioactif, donne en se désintégrant un noyau d'argon
          $\ce{_{18}^{36}Ar}$.
          \begin{enumerate}
              \item Écrire l'équation de désintégration d'un noyau de chlore
                    36.~\textbf{(1pt)}
              \item Identifier le type de cette désintégration.~\textbf{(0,5pt)}
              \item Calculer, en unité (\unit{\MeV}), l'énergie libérée
                    $E_{\text{libérée}}=|\Delta E|$ au cours de la
                    désintégration d'un noyau de chlore 36.~\textbf{(1pt)}
          \end{enumerate}
\end{enumerate}

